% Source: https://artofproblemsolving.com/community/c3h2518698_raised_empty_dotor_multiplication_symbol
% Stable identity: aops-topic-2518698-q01
\textbf{(a)} $G(2)=4-3=1$, so $(F\circ G)(2)=F(1)=3$. $F(2)=5$, so $(G\circ F)(2)=G(5)=25-3=22$.

\textbf{(b)}
\[(F\circ G)(x)=2\left(x^2-3\right)+1=2x^2-5,\qquad(G\circ F)(x)=(2x+1)^2-3=4x^2+4x-2.\]

\textbf{(c)} Equality would mean $2x^2-5=4x^2+4x-2$, i.e. $2x^2+4x+3=0$. Its discriminant is $16-24=-8<0$; equivalently $2x^2+4x+3=2(x+1)^2+1\ge1>0$. So there is no real solution: here the two compositions never agree, and in general $F\circ G\ne G\circ F$ (composition is not commutative).

$\boxed{(F\circ G)(2)=3,\ (G\circ F)(2)=22;\ F\circ G=2x^2-5,\ G\circ F=4x^2+4x-2;\ \text{never equal}}$
